\documentclass[12pt]{article}
% Usar dentro do SBC Conferences Template oficial do Overleaf (sbc-template.sty e sbc.bst).
\usepackage{sbc-template}
\usepackage[utf8]{inputenc}
\usepackage[T1]{fontenc}
\usepackage[brazil]{babel}
\usepackage{graphicx,url,booktabs,array,amsmath}
\usepackage{tikz,pgfplots}
\pgfplotsset{compat=1.18}
\usepackage[hidelinks]{hyperref}
\sloppy
% Preencher somente com URL verdadeiro e acessível, vídeo publicado e link validado.
\newcommand{\URLVideo}{https://www.youtube.com/watch?v=GYhiJRu9faM}

\newcommand{\VideoNaoPublicado}{}
\title{RANGOT: classificação de escopo e intenção com\\recuperação documental para o restaurante universitário}
\author{Cleiton Pinheiro Rodrigues\inst{1}, Jonathan Pereira do Amaral\inst{1}\\
Felippe Baes Sales Lara\inst{1}, Gabriel Mattos de Aquino\inst{1}}
\address{Engenharia de Software -- FACOM -- Faculdade de Computação\\
Universidade Federal de Mato Grosso do Sul (UFMS)\\
RGAs, na ordem de autoria: 202419060102; 202419060072;\\
202219060141; 202419060064}
\begin{document}
\maketitle
\begin{abstract}
RANGOT is an academic chatbot about the university restaurant at UFMS.
A scope classifier precedes a six-intent classifier; LangGraph controls routing
and document retrieval supports language-model responses. Training comprises
550 messages, including 268 human-written ones; the held-out test contains 160
human messages. Five-fold cross-validation compares three algorithms and two
representations. Calibrated character-based SVMs achieve macro-F1 scores of
0.9229 and 0.9200. A single independent evaluation yields end-to-end macro-F1
of 0.9499, 149 correct routes and seven clarification requests. We discuss
the small numerical advantage of an eight-class model in training CV,
routing errors, document authority, synthetic provenance and selection bias.
\end{abstract}
\begin{resumo}
O RANGOT é um chatbot acadêmico sobre o Restaurante Universitário da UFMS.
Um classificador de escopo antecede outro de seis intenções; LangGraph controla
as rotas e a recuperação documental fundamenta respostas do modelo de linguagem.
O treino contém 550 mensagens, 268 humanas, e o teste reúne 160 mensagens humanas.
Comparam-se três algoritmos e duas representações em validação cruzada de cinco
partições. SVMs calibradas com caracteres alcançam F1 macro de 0,9229 e 0,9200.
Na avaliação independente única, o grafo obtém F1 macro 0,9499, 149 rotas corretas
e sete esclarecimentos. Discutem-se a pequena vantagem numérica de um modelo
de oito classes na CV, erros, autoridade documental, proveniência e viés de seleção.
\end{resumo}

\section{Introdução}
Dúvidas sobre cardápio, horários, subsídio e acesso ao Restaurante Universitário
(RU) exigem identificar o assunto e consultar fontes com públicos e vigências
distintos. Classificação errada ou documentação antiga podem produzir orientações
inadequadas. O RANGOT, protótipo acadêmico de Inteligência Artificial da
FACOM/UFMS em 2026/2, investiga esses problemas.

Suas contribuições são uma base curada com proveniência explícita, comparação
de classificadores, estudo de cascata versus modelo único e política empírica de
esclarecimento. Os classificadores determinam a rota, o LangGraph coordena a
execução e o RAG fornece evidências ao modelo de linguagem (LLM).

\section{Fundamentação}
Na classificação supervisionada, aprende-se uma função entre textos e rótulos.
TF-IDF pondera a frequência dos termos pela raridade na coleção; palavras e
bigramas representam unidades lexicais, enquanto caracteres compartilham
subsequências de abreviações e grafias variantes \cite{manning2008}.
Isso não garante compreensão contextual.

A regressão logística aprende fronteiras por perda probabilística regularizada.
Complement Naive Bayes estima pesos pelo complemento das classes
\cite{rennie2003}; a SVM aprende separadores de margem \cite{cortes1995}.
Como \texttt{LinearSVC} não fornece probabilidades, utiliza-se calibração
sigmoide com partições próprias \cite{sklearncalibration} em pipelines
scikit-learn \cite{pedregosa2011}.

Acurácia é a proporção de acertos. F1 combina precisão e revocação pela média
harmônica; sua média macro atribui igual peso às classes. A validação cruzada
(CV) estratificada preserva aproximadamente suas proporções, mas selecionar
e descrever o vencedor com as mesmas partições introduz viés \cite{varma2006}.

Na cascata, rejeições de escopo impedem o estágio de intenção.
LangGraph organiza estado e transições condicionais \cite{langgraph}.
RAG condiciona a geração a material recuperado \cite{lewis2020}.
Aqui, TF-IDF e cosseno recuperam trechos após o roteamento; o LLM gera texto
condicionado ao contexto. Fluência e recuperação não comprovam fidelidade factual.

\section{Domínio, dados e documentos}
As intenções são \texttt{cardapio}, \texttt{horarios}, \texttt{valores},
\texttt{recarga\_creditos}, \texttt{cadastro} e \texttt{regras\_uso}.
Elas formam o escopo \texttt{dominio}; \texttt{chitchat} representa conversa
social e \texttt{fora\_dominio}, solicitações externas. O rótulo único limita
o tratamento de perguntas compostas.

O Google Forms recebeu 67 envios e 669 respostas preenchidas. A curadoria
reservou 31 envios para teste e 36 para treino, com quarentena de sobreposições,
antes dos experimentos. O treino tem 268 textos humanos e 282 de LLM (48,73\%
humanos); o teste conserva 160 textos humanos originais (Tabela~\ref{tab:dados}).
\begin{table}[!htb]
\centering
\caption{Composição dos conjuntos congelados. Teste: composição registrada no protocolo de curadoria; nenhuma frase sintética.}
\label{tab:dados}
\begin{tabular}{lrrrr}
\toprule Intenção & Treino & Humano & LLM & Teste\\ \midrule
cardapio & 50 & 40 & 10 & 15\\
horarios & 50 & 34 & 16 & 15\\
valores & 50 & 32 & 18 & 15\\
recarga\_creditos & 50 & 32 & 18 & 15\\
cadastro & 50 & 26 & 24 & 15\\
regras\_uso & 50 & 26 & 24 & 15\\
chitchat & 100 & 25 & 75 & 30\\
fora\_dominio & 150 & 53 & 97 & 40\\
\midrule Total & 550 & 268 & 282 & 160\\ \bottomrule
\end{tabular}
\end{table}


Revisaram-se rótulos, ambiguidades, múltiplos pedidos, duplicatas exatas,
normalizadas e sinais lexicais de paráfrases. Decisões delegadas ao Codex
foram conferidas e ratificadas; treze pendências não foram aprovadas
e uma candidata semelhante foi excluída. Revisão humana não altera autoria:
frases sintéticas mantêm \texttt{origem=llm}.

Negativos difíceis incluem biblioteca, secretaria, bolsas, matrícula e transporte,
com vocabulário próximo ao RU. A síntese complementar
recebeu somente treino autorizado em conversa separada; não há comprovação de
isolamento físico. Ferramentas e prompts de 158 frases LLM (157 originais e uma
posterior) têm registro incompleto, limitando a reprodução da geração.

A base ativa tem dez documentos, oito sínteses/orientações e dois PDFs, divididos
em 83 trechos e catalogados em \texttt{documentos/FONTES.md}.
O Guia Prático \cite{ufmsguia} fundamenta horários, elegibilidade e serviço;
a Resolução 399/2025 trata de assistência estudantil \cite{ufms399}, e o
Edital 114/2026 atende ingressantes internacionais \cite{ufms114}.
Pix e cartões são relatos atribuídos ao grupo, não normas. Sem cardápio vigente
para data e campus, informa-se o canal oficial da PROAES.

\section{Metodologia}
\subsection{Classificadores e seleção}
O escopo usa 550 mensagens; o tema, somente as 300 originalmente do domínio.
Comparam-se logística, Complement NB e SVM calibrada, com TF-IDF de palavras
$(1,2)$ ou caracteres \texttt{char\_wb} $(2,5)$. O pré-processamento usa
minúsculas, remove acentos/pontuação e regulariza espaços; preserva stopwords,
adota frequência sublinear e frequência documental mínima um. A normalização
também elimina pistas de pontuação.

Cinco partições estratificadas, embaralhadas com semente 42, são reutilizadas
entre candidatos. GridSearchCV compara \(C\in\{0{,}1;1;10\}\) na logística/SVM
e \(\alpha\in\{0{,}1;0{,}5;1\}\) no NB: 18 configurações por estágio.
Seleciona-se F1 macro; desempates usam acurácia, menor desvio e ordem estável.
Logística/SVM têm pesos balanceados. Médias e desvios populacionais entre
partições não são intervalos de confiança.

A calibração sigmoide envolve todo o pipeline vetorizador/SVM, com três
partições internas e \texttt{ensemble=False}, sem ajustar vocabulário, IDF ou
calibrador na validação. Matrizes e análises por origem usam previsões
fora de partição (OOF) do vencedor após a busca, sujeitas ao viés de seleção.
Os modelos finais usam todos os exemplos pertinentes.

\subsection{A1, confiança e integração}
O A1 compara cascata e modelo único de oito classes nas mesmas cinco partições
externas estratificadas. Cada treino externo seleciona \(C\) em três partições;
a calibração tem três partições próprias. O tema usa apenas domínio verdadeiro.
Família e representação foram escolhidas anteriormente nesse conjunto:
a busca interna não elimina todo o viés.

Outra execução OOF usa a cascata de produção (\(C=1\) no escopo;
\(C=0{,}1\) no tema). Entre 18 limiares, minimiza-se custo: falso aceite 5,
intenção errada 3, esclarecimento 1, falsa recusa 2 e troca social/externo 1.
Pesos foram declarados antes da execução, sem validação com usuários;
empates favorecem acertos diretos e menor limiar. Seleção e descrição usam
os mesmos OOF.

O grafo solicita esclarecimento quando a confiança do tema é inferior a
\(0{,}60\). O RAG filtra intenção e documentos gerais, usa caracteres $(3,5)$,
trechos de cerca de 700 caracteres, até quatro resultados e cosseno mínimo 0,05.
Fontes acompanham o contexto, separado das instruções do sistema.
Sem LLM há respostas fixas; sem evidência, orienta-se admitir a lacuna.
Não há consulta online durante a conversa.

\subsection{Avaliação e reprodutibilidade}
A avaliação oficial única foi executada manualmente, sem LLM, após congelar
modelos, configuração e teste. Mede escopo nas 160 mensagens, tema-oráculo nas
90 realmente do domínio e rota completa, além de falsas aceitações/recusas,
esclarecimento e latência. JSON/CSV foram preservados, sem ajuste posterior.
O teste local não é a união de conjuntos do professor para C7.

Na reprodução manual, as 22 células executaram em novo kernel e cópia isolada com Python 3.12.13/scikit-learn 1.9.1. Três linhas do A1 tiveram diferença máxima de 0,0, e 39 arquivos protegidos permaneceram intactos. Os modelos reproduzidos foram validados sem substituir os oficiais. As suítes na cópia e no checkout aprovaram 160 testes cada.


\section{Resultados}
\subsection{Escopo e tema}
SVM/caracteres venceu em ambos os estágios: F1 macro 0,922867 no escopo e
0,919959 no tema, com \(C=1\) e \(C=0{,}1\), respectivamente
(Tabela~\ref{tab:comparacao}). A diferença para logística/caracteres no tema é
apenas 0,000416; o ranking não comprova superioridade estatística.
Caracteres beneficiaram logística e SVM no escopo, mas não Complement NB:
o efeito depende da combinação de algoritmo e representação.
\begin{table}[!htb]
\centering\small
\caption{CV interna: médias e desvios populacionais de cinco partições. \(p\): \(C\) na logística/SVM; \(\alpha\) no Complement NB.}
\label{tab:comparacao}
\begin{tabular}{lllccc}
\toprule Estágio & Algoritmo & TF-IDF & F1 macro & Acurácia & \(p\)\\ \midrule
Escopo & SVM cal. & Carac. & \(0{,}9229\pm0{,}0133\) & \(0{,}9309\pm0{,}0148\) & 1,0\\
Escopo & Logística & Carac. & \(0{,}9150\pm0{,}0173\) & \(0{,}9236\pm0{,}0159\) & 10,0\\
Escopo & SVM cal. & Palav. & \(0{,}8977\pm0{,}0192\) & \(0{,}9055\pm0{,}0196\) & 10,0\\
Escopo & Logística & Palav. & \(0{,}8664\pm0{,}0356\) & \(0{,}8855\pm0{,}0285\) & 10,0\\
Escopo & Comp. NB & Palav. & \(0{,}8491\pm0{,}0458\) & \(0{,}8673\pm0{,}0417\) & 0,1\\
Escopo & Comp. NB & Carac. & \(0{,}8486\pm0{,}0206\) & \(0{,}8655\pm0{,}0194\) & 0,1\\
\midrule
Tema & SVM cal. & Carac. & \(0{,}9200\pm0{,}0245\) & \(0{,}9200\pm0{,}0245\) & 0,1\\
Tema & Logística & Carac. & \(0{,}9195\pm0{,}0189\) & \(0{,}9200\pm0{,}0194\) & 10,0\\
Tema & Comp. NB & Carac. & \(0{,}9093\pm0{,}0221\) & \(0{,}9100\pm0{,}0226\) & 0,5\\
Tema & Comp. NB & Palav. & \(0{,}8980\pm0{,}0134\) & \(0{,}9000\pm0{,}0149\) & 0,1\\
Tema & SVM cal. & Palav. & \(0{,}8837\pm0{,}0183\) & \(0{,}8867\pm0{,}0163\) & 10,0\\
Tema & Logística & Palav. & \(0{,}8807\pm0{,}0234\) & \(0{,}8833\pm0{,}0211\) & 10,0\\
\bottomrule\end{tabular}
\end{table}


\begin{figure}[!htb]
\centering
% Eixos previstos abaixo da grade; títulos alinhados no topo dos painéis.
\begin{minipage}[t]{.48\linewidth}\centering
\textbf{Escopo (550 OOF)}\par
\begin{tikzpicture}[x=.90cm,y=.90cm,font=\footnotesize]
\filldraw[fill=blue!85,draw=white] (0,2) rectangle (1,3);
\node[text=white] at (0.5,2.5) {285};
\filldraw[fill=blue!2,draw=white] (1,2) rectangle (2,3);
\node[text=black] at (1.5,2.5) {6};
\filldraw[fill=blue!3,draw=white] (2,2) rectangle (3,3);
\node[text=black] at (2.5,2.5) {9};
\filldraw[fill=blue!2,draw=white] (0,1) rectangle (1,2);
\node[text=black] at (0.5,1.5) {6};
\filldraw[fill=blue!27,draw=white] (1,1) rectangle (2,2);
\node[text=black] at (1.5,1.5) {90};
\filldraw[fill=blue!1,draw=white] (2,1) rectangle (3,2);
\node[text=black] at (2.5,1.5) {4};
\filldraw[fill=blue!3,draw=white] (0,0) rectangle (1,1);
\node[text=black] at (0.5,0.5) {11};
\filldraw[fill=blue!1,draw=white] (1,0) rectangle (2,1);
\node[text=black] at (1.5,0.5) {2};
\filldraw[fill=blue!41,draw=white] (2,0) rectangle (3,1);
\node[text=black] at (2.5,0.5) {137};
\node[anchor=north] at (0.5,-.15) {Dom};
\node[anchor=north] at (1.5,-.15) {Soc};
\node[anchor=north] at (2.5,-.15) {Ext};
\node[anchor=east] at (-.15,2.5) {Dom};
\node[anchor=east] at (-.15,1.5) {Soc};
\node[anchor=east] at (-.15,0.5) {Ext};
\node at (1.5,-.95) {Previsto};
\node[rotate=90] at (-1.55,1.5) {Real};
\end{tikzpicture}\end{minipage}
\hfill
\begin{minipage}[t]{.48\linewidth}\centering
\textbf{Tema (300 OOF)}\par
\begin{tikzpicture}[x=.72cm,y=.72cm,font=\footnotesize]
\filldraw[fill=blue!81,draw=white] (0,5) rectangle (1,6);
\node[text=white] at (0.5,5.5) {46};
\filldraw[fill=blue!2,draw=white] (1,5) rectangle (2,6);
\node[text=black] at (1.5,5.5) {1};
\filldraw[fill=blue!2,draw=white] (2,5) rectangle (3,6);
\node[text=black] at (2.5,5.5) {1};
\filldraw[fill=blue!0,draw=white] (3,5) rectangle (4,6);
\node[text=black] at (3.5,5.5) {0};
\filldraw[fill=blue!2,draw=white] (4,5) rectangle (5,6);
\node[text=black] at (4.5,5.5) {1};
\filldraw[fill=blue!2,draw=white] (5,5) rectangle (6,6);
\node[text=black] at (5.5,5.5) {1};
\filldraw[fill=blue!5,draw=white] (0,4) rectangle (1,5);
\node[text=black] at (0.5,4.5) {3};
\filldraw[fill=blue!81,draw=white] (1,4) rectangle (2,5);
\node[text=white] at (1.5,4.5) {46};
\filldraw[fill=blue!2,draw=white] (2,4) rectangle (3,5);
\node[text=black] at (2.5,4.5) {1};
\filldraw[fill=blue!0,draw=white] (3,4) rectangle (4,5);
\node[text=black] at (3.5,4.5) {0};
\filldraw[fill=blue!0,draw=white] (4,4) rectangle (5,5);
\node[text=black] at (4.5,4.5) {0};
\filldraw[fill=blue!0,draw=white] (5,4) rectangle (6,5);
\node[text=black] at (5.5,4.5) {0};
\filldraw[fill=blue!5,draw=white] (0,3) rectangle (1,4);
\node[text=black] at (0.5,3.5) {3};
\filldraw[fill=blue!2,draw=white] (1,3) rectangle (2,4);
\node[text=black] at (1.5,3.5) {1};
\filldraw[fill=blue!80,draw=white] (2,3) rectangle (3,4);
\node[text=white] at (2.5,3.5) {45};
\filldraw[fill=blue!0,draw=white] (3,3) rectangle (4,4);
\node[text=black] at (3.5,3.5) {0};
\filldraw[fill=blue!0,draw=white] (4,3) rectangle (5,4);
\node[text=black] at (4.5,3.5) {0};
\filldraw[fill=blue!2,draw=white] (5,3) rectangle (6,4);
\node[text=black] at (5.5,3.5) {1};
\filldraw[fill=blue!2,draw=white] (0,2) rectangle (1,3);
\node[text=black] at (0.5,2.5) {1};
\filldraw[fill=blue!0,draw=white] (1,2) rectangle (2,3);
\node[text=black] at (1.5,2.5) {0};
\filldraw[fill=blue!4,draw=white] (2,2) rectangle (3,3);
\node[text=black] at (2.5,2.5) {2};
\filldraw[fill=blue!80,draw=white] (3,2) rectangle (4,3);
\node[text=white] at (3.5,2.5) {45};
\filldraw[fill=blue!2,draw=white] (4,2) rectangle (5,3);
\node[text=black] at (4.5,2.5) {1};
\filldraw[fill=blue!2,draw=white] (5,2) rectangle (6,3);
\node[text=black] at (5.5,2.5) {1};
\filldraw[fill=blue!0,draw=white] (0,1) rectangle (1,2);
\node[text=black] at (0.5,1.5) {0};
\filldraw[fill=blue!0,draw=white] (1,1) rectangle (2,2);
\node[text=black] at (1.5,1.5) {0};
\filldraw[fill=blue!2,draw=white] (2,1) rectangle (3,2);
\node[text=black] at (2.5,1.5) {1};
\filldraw[fill=blue!4,draw=white] (3,1) rectangle (4,2);
\node[text=black] at (3.5,1.5) {2};
\filldraw[fill=blue!81,draw=white] (4,1) rectangle (5,2);
\node[text=white] at (4.5,1.5) {46};
\filldraw[fill=blue!2,draw=white] (5,1) rectangle (6,2);
\node[text=black] at (5.5,1.5) {1};
\filldraw[fill=blue!2,draw=white] (0,0) rectangle (1,1);
\node[text=black] at (0.5,0.5) {1};
\filldraw[fill=blue!0,draw=white] (1,0) rectangle (2,1);
\node[text=black] at (1.5,0.5) {0};
\filldraw[fill=blue!2,draw=white] (2,0) rectangle (3,1);
\node[text=black] at (2.5,0.5) {1};
\filldraw[fill=blue!0,draw=white] (3,0) rectangle (4,1);
\node[text=black] at (3.5,0.5) {0};
\filldraw[fill=blue!0,draw=white] (4,0) rectangle (5,1);
\node[text=black] at (4.5,0.5) {0};
\filldraw[fill=blue!85,draw=white] (5,0) rectangle (6,1);
\node[text=white] at (5.5,0.5) {48};
\node[anchor=north] at (0.5,-.15) {Car};
\node[anchor=north] at (1.5,-.15) {Hor};
\node[anchor=north] at (2.5,-.15) {Val};
\node[anchor=north] at (3.5,-.15) {Rec};
\node[anchor=north] at (4.5,-.15) {Cad};
\node[anchor=north] at (5.5,-.15) {Reg};
\node[anchor=east] at (-.15,5.5) {Car};
\node[anchor=east] at (-.15,4.5) {Hor};
\node[anchor=east] at (-.15,3.5) {Val};
\node[anchor=east] at (-.15,2.5) {Rec};
\node[anchor=east] at (-.15,1.5) {Cad};
\node[anchor=east] at (-.15,0.5) {Reg};
\node at (3.0,-.95) {Previsto};
\node[rotate=90] at (-1.55,3.0) {Real};
\end{tikzpicture}\end{minipage}

\caption{Contagens OOF pós-seleção: linhas reais, colunas previstas.
Dom/Soc/Ext: domínio/social/externo; Car/Hor/Val/Rec/Cad/Reg: cardápio,
horários, valores, recarga, cadastro e regras. Cores relativas ao máximo de cada matriz.}
\label{fig:matrizes}
\end{figure}
Na Figura~\ref{fig:matrizes}, o escopo cometeu 38 erros, incluindo 17 falsos aceites em 250 não-domínio
(6,8\%) e 15 falsas recusas em 300 domínio (5,0\%). O tema teve 24 erros;
horários e valores foram confundidos com cardápio três vezes cada.
Cardápio teve F1 0,884615 e regras, 0,941176. Matrizes agregadas e médias
por partição não são a mesma agregação.

\begin{figure}[!htb]
\centering
\begin{minipage}{.49\linewidth}\centering
\begin{tikzpicture}
\begin{axis}[width=\linewidth,height=4.2cm,font=\footnotesize,xmode=log,log basis x=10,
xlabel={\(C\)},ylabel={F1 macro},xtick={0.1,1,10},xticklabels={$0{,}1$,$1$,$10$},
ymin=.8,ymax=1.02,grid=major,legend style={font=\footnotesize,at={(.98,.04)},anchor=south east},
title={Escopo}]
\addplot+[black,mark=square*,error bars/.cd,y dir=both,y explicit] coordinates { (0.1,0.9822473061345148) +- (0,0.003634587148130988) (1.0,1) +- (0,0) (10.0,1) +- (0,0) };
\addlegendentry{Treino}
\addplot+[blue,mark=*,error bars/.cd,y dir=both,y explicit] coordinates { (0.1,0.9122824153411486) +- (0,0.009317981022236603) (1.0,0.9228672897685202) +- (0,0.01329960679823014) (10.0,0.9156991350222699) +- (0,0.01374917251325642) };
\addlegendentry{Validação}
\end{axis}\end{tikzpicture}\end{minipage}
\hfill
\begin{minipage}{.49\linewidth}\centering
\begin{tikzpicture}
\begin{axis}[width=\linewidth,height=4.2cm,font=\footnotesize,xmode=log,log basis x=10,
xlabel={\(C\)},ylabel={F1 macro},xtick={0.1,1,10},xticklabels={$0{,}1$,$1$,$10$},
ymin=.8,ymax=1.02,grid=major,legend style={font=\footnotesize,at={(.98,.04)},anchor=south east},
title={Tema}]
\addplot+[black,mark=square*,error bars/.cd,y dir=both,y explicit] coordinates { (0.1,0.9933326821899253) +- (0,0.002041241452319363) (1.0,1) +- (0,0) (10.0,1) +- (0,0) };
\addlegendentry{Treino}
\addplot+[blue,mark=*,error bars/.cd,y dir=both,y explicit] coordinates { (0.1,0.9199591534831809) +- (0,0.02450762891555043) (1.0,0.9195886477929822) +- (0,0.02666491561130133) (10.0,0.9159272214999771) +- (0,0.01792968772142712) };
\addlegendentry{Validação}
\end{axis}\end{tikzpicture}\end{minipage}

\caption{SVM/caracteres: F1 médio de treino e validação, com desvio populacional
de cinco partições; eixo vertical de 0,80 a 1,02.}
\label{fig:curvas}
\end{figure}
Aumentar \(C\) leva F1 de treino a 1,0 sem melhoria monotônica da validação
(Figura~\ref{fig:curvas}). No tema, \(C=0{,}1\) obtém 0,919959;
\(C=10\), 0,915927. A separação treino/validação é compatível com sobreajuste,
sem explicar causalmente cada erro.

\subsection{Cascata, modelo único e limiar}
\begin{table}[!htb]
\centering\small
\caption{A1 sem limiar: médias e desvios entre partições; contagens de erros OOF. FA: falso aceite entre 250 não-domínio; FR: falsa recusa entre 300 domínio. A referência 08 usa \(C\) fixos.}
\label{tab:a1}
\begin{tabular}{lccrrr}
\toprule Arquitetura & F1 macro & Acurácia & Erros & FA & FR\\ \midrule
Cascata, busca interna & \(0{,}8821\pm0{,}0140\) & \(0{,}8891\pm0{,}0106\) & 61 & 20 & 13\\
Único, busca interna & \(0{,}8851\pm0{,}0189\) & \(0{,}8909\pm0{,}0191\) & 60 & 17 & 20\\
Referência 08 & \(0{,}8787\pm0{,}0096\) & \(0{,}8855\pm0{,}0093\) & 63 & 19 & 16\\
\bottomrule\end{tabular}
\end{table}

No A1 (Tabela~\ref{tab:a1}), os F1 médios foram 0,882105 (cascata) e 0,885115 (único):
vantagem de 0,003010 e um erro OOF a menos. O único venceu quatro partições,
mas perdeu cerca de 0,0427 na terceira (Figura~\ref{fig:a1}).
A cascata perdeu 13 perguntas no escopo; o tema isolado acertaria 11 delas.
Outras 22 passaram e tiveram intenção errada. O único reduziu falsos aceites
de 20 para 17 e elevou falsas recusas de 13 para 20, sem teste de significância.

O único exige um pipeline; a cascata, dois. Registros históricos indicam
ajuste de 1,246 s e 0,070 ms/mensagem batch na cascata, contra 0,920 s e
0,037 ms no único, cujo tamanho serializado em memória foi ligeiramente maior.
São diagnósticos amostrais; executar o tema em todo o batch não mede latência
da rota online nem da API.

\begin{figure}[!htb]
\centering
\begin{minipage}{.49\linewidth}\centering
\begin{tikzpicture}\begin{axis}[width=\linewidth,height=4.2cm,font=\footnotesize,
xlabel={Partição externa},ylabel={F1 macro},xmin=.8,xmax=5.2,xtick={1,2,3,4,5},ymin=.8,ymax=.93,grid=major,
legend style={font=\footnotesize,at={(.98,.04)},anchor=south east}]
\addplot+[black,mark=square*] coordinates { (1,0.8976980590180489) (2,0.8901018704222365) (3,0.8911027568922306) (4,0.8713707902267225) (5,0.8602518364128271) };
\addlegendentry{Cascata}
\addplot+[blue,mark=*] coordinates { (1,0.8997940388842228) (2,0.8991153127987476) (3,0.8484414968298989) (4,0.8893993705193215) (5,0.8888257992296718) };
\addlegendentry{Único}
\end{axis}\end{tikzpicture}\end{minipage}\hfill
\begin{minipage}{.49\linewidth}\centering
\begin{tikzpicture}\begin{axis}[width=\linewidth,height=4.2cm,font=\footnotesize,
xlabel={Limiar},ylabel={Mensagens},xmin=0,xmax=1,ymin=0,ymax=320,grid=major,
legend style={font=\footnotesize,at={(.02,.98)},anchor=north west}]
\addplot+[black,mark=square*,mark size=1pt] coordinates { (0,41) (0.2,41) (0.25,41) (0.3,41) (0.35,40) (0.4,40) (0.45,37) (0.5,32) (0.55,25) (0.6,21) (0.65,19) (0.7,18) (0.75,16) (0.8,13) (0.85,9) (0.9,6) (0.95,3) (1,0) };
\addlegendentry{Erros diretos}
\addplot+[blue,mark=*,mark size=1pt] coordinates { (0,0) (0.2,0) (0.25,0) (0.3,0) (0.35,1) (0.4,1) (0.45,6) (0.5,14) (0.55,24) (0.6,32) (0.65,46) (0.7,55) (0.75,68) (0.8,86) (0.85,102) (0.9,141) (0.95,210) (1,303) };
\addlegendentry{Esclarecimentos}
\addplot[dashed,gray,forget plot] coordinates {(.6,0) (.6,320)};
\end{axis}\end{tikzpicture}\end{minipage}

\caption{A1 sem limiar: F1 por partição externa (esquerda).
Referência 08: erros diretos e esclarecimentos OOF (direita);
linha vertical em \(0{,}60\). Escalas verticais distintas.}
\label{fig:a1}
\end{figure}
Na referência 08, o limiar reduziu respostas diretas de 303 para 271 e erros
diretos de 41 para 21: 32 esclarecimentos evitaram 20 erros e bloquearam
12 acertos. A acurácia condicional subiu de 86,47\% para 92,25\%;
o custo caiu de 199 para 159. Porém, a acurácia completa caiu de 0,885455
para 0,863636; F1 agregado passou de 0,878652 para 0,883098.
Troca-se cobertura por cautela, sem melhorar todas as métricas.
Custos maiores para erros diretos selecionariam 0,85 na análise de sensibilidade.
A produção mantém os dois modelos oficiais.

\subsection{Avaliação independente}
O grafo alcançou F1 macro \(0{,}949881\) e acurácia \(0{,}931250\)
(Tabela~\ref{tab:independente}).
\newcommand{\ConclusaoTeste}{No teste humano local, houve 149 acertos em 160 mensagens, F1 macro de \(0{,}9499\) e sete esclarecimentos.}
\begin{table}[!htb]\centering\small
\caption{CV/OOF de treino e avaliação independente local. \(N\) refere-se ao teste. Ponta a ponta: OOF agregado da referência 08 com limiar; não é a média do A1 com busca interna.}
\label{tab:independente}
\begin{tabular}{lrcccc}\toprule
Nível & \(N\) & F1 CV/OOF & Acurácia CV/OOF & F1 teste & Acurácia teste\\ \midrule
Escopo & 160 & \(0{,}9229\pm0{,}0133\) & \(0{,}9309\pm0{,}0148\) & \(0{,}9633\) & \(0{,}9688\)\\
Tema-oráculo & 90 & \(0{,}9200\pm0{,}0245\) & \(0{,}9200\pm0{,}0245\) & \(0{,}9667\) & \(0{,}9667\)\\
Ponta a ponta, \(0{,}60\) & 160 & \(0{,}8831\) & \(0{,}8636\) & \(0{,}9499\) & \(0{,}9313\)\\
\bottomrule\end{tabular}\end{table}
No teste reservado, os classificadores isolados acertaram 155/160 mensagens de escopo e 87/90 de tema-oráculo. O grafo completo acertou 149/160. As taxas finais foram \(1/70=0{,}014286\) de falso aceite, \(1/90=0{,}011111\) de falsa recusa e \(7/160=0{,}04375\) de esclarecimento. A latência média registrada foi \(1{,}244327\) ms por mensagem, sem LLM; não é latência de API nem comparável diretamente ao batch do A1.
O F1 macro oficial é a média nas oito classes reais. A rota \texttt{esclarecer} reduz a revocação da classe verdadeira, mas não entra como nona classe na média. As ponderações distintas permitem F1 maior que acurácia. Esclarecimento continua sendo divergência quando se espera uma classe específica.


\section{Discussão e ameaças à validade}
\subsection{Erros e generalização}
F1 de escopo e tema-oráculo no teste superaram a CV em 0,040426 e 0,046695.
O ganho ponta a ponta de 0,066783 é relativo ao OOF agregado da referência 08
com limiar, não ao estimador A1 com busca interna. O modelo único não foi
avaliado no teste; o oráculo de 90 mensagens não representa o fluxo completo.
Protocolos e distribuições distintos impedem inferir ganho geral, ausência de
sobreajuste ou significância estatística.

Dos 11 erros, cinco são esclarecimentos de domínio (dois cardápio, dois valores,
um cadastro). Três tinham tema correto, com confiança entre 0,5699 e 0,5896;
dois bloquearam tema errado. Os demais são horários$\to$cardápio,
horários$\to$social, externo$\to$horários, externo$\to$social e dois
esclarecimentos não-domínio (externo e social). Esses grupos são disjuntos.
Os sete esclarecimentos podem ser conservadores na aplicação, mas divergem
dos rótulos esperados no protocolo.

“Que q vai ter no RU hj?” teve tema correto e esclarecimento (0,5699).
“tem ru hoje?” recebeu cardápio contra horários (0,9613), e horas complementares
recebeu horários (0,9045): confiança alta não impede erro.
“Boia” apareceu na falsa recusa; “Bora comer no RU hj?” era social e foi
esclarecida. Abreviações, elipses e sobreposição lexical são hipóteses,
SEM atribuição causal. Rótulos, previsões e limiar não foram ajustados.

Nos OOF de tema, humanos tiveram F1 0,939947 (\(n=190\)) e sintéticos
0,884631 (\(n=110\)); no escopo, 0,927125 e 0,916766.
Comprimentos médios foram 48,72 e 43,29 caracteres. Classe, estilo,
dificuldade e seleção confundem a comparação: não houve testes isolados de origem.

Coleta por conveniência e curadoria não representam frequências de uso.
Separar envios reduz compartilhamento, mas não comprova autores únicos.
Autores e famílias sintéticas não foram agrupados na CV; relações entre frases
podem elevar escores, e a triagem lexical não exclui toda paráfrase.
Reprodução exata demonstra consistência de execução, não ausência desses vieses.
Menor dificuldade do teste é hipótese; não se extrapolam resultados a conjuntos
ocultos. Falsos aceites arriscam orientações inadequadas, recusas reduzem utilidade
e esclarecimentos excessivos oneram a interação. Custos não medem impacto real
sobre estudantes; a arquitetura continua limitada a um rótulo por mensagem.

\subsection{Integração e limites documentais}
Nos dez cenários fictícios históricos do A1, houve cinco rotas esperadas entre
oito casos de classe única: cardápio/pagamento foram para valores; cadastro,
para esclarecimento. O caso vago foi social e o composto recebeu uma intenção.
A Missão 10 usou componentes controlados e nove perguntas com modelos reais
e LLM simulado. Esses registros não constituem taxa padronizada da API.
O registro antigo de oito documentos/77 trechos antecede a base atual de dez/83.

Os testes foram confirmados de forma manual com GPT-6 Luna/OpenRouter e sem LLM,
sem estudo sistemático da geração. Recuperação pode não responder à pergunta,
e intenção errada restringe fontes. Prompts e atribuição não garantem ausência
de alucinações ou resistência universal a instruções maliciosas.

Persistem lacunas sobre pagamento oficial, preço integral em Campo Grande,
condutas e horários conflitantes de outros campi. Cardápio planejado não prova
refeição servida. O chatbot não solicita senhas/documentos pessoais nem executa
cadastros. Avaliar recuperação anotada e fidelidade com critérios humanos é
trabalho futuro, não resultado do A1.

\section{Conclusão}
O RANGOT combina classificação em cascata, esclarecimento e recuperação documental.
SVM/caracteres venceu a comparação; o A1 revelou pequena diferença numérica
e propagação de erros. No OOF, o limiar reduziu erros diretos e cobertura.
No teste humano local, houve 149 acertos em 160 mensagens,
F1 macro de 0,9499 e sete esclarecimentos.
A reprodução preservou os artefatos oficiais. Limitações de coleta, seleção
e proveniência motivam estudos com controle de autorias/famílias, custos com
usuários e fidelidade documental.
\footnote{Vídeo: \href{https://www.youtube.com/watch?v=GYhiJRu9faM}{demonstração do RANGOT}.}

\subsection*{Contribuição dos integrantes}
\textbf{Cleiton Pinheiro Rodrigues (202419060102)} atuou em desenvolvimento,
integração, classificadores, experimentos, A1, testes, curadoria, documentação,
revisão humana e reprodução.
\textbf{Jonathan Pereira do Amaral (202419060072)} atuou em desenvolvimento,
organização e manutenção do GitHub/repositório, acompanhamento de alterações,
análise de pull requests, revisão e aprovação de contribuições, além de auditoria
independente da reprodução.
\textbf{Felippe Baes Sales Lara (202219060141)} e
\textbf{Gabriel Mattos de Aquino (202419060064)} atuaram na elaboração e revisão
do Google Forms, divulgação da coleta e apoio à obtenção das
respostas. Também participaram das discussões e decisões gerais do projeto,
da análise de pull requests, da elaboração e revisão do artigo e da avaliação
de propostas e alternativas durante o desenvolvimento do trabalho.

\section*{Declaração de uso de IA generativa}
Cleiton usou ChatGPT/GPT-6 e Codex/GPT-6.1-Sol, raciocínio High em planejamento, requisitos, código, configurações,
curadoria, experimentos, análise, documentação e reprodução.
O \texttt{xhigh} persistido não comprova o modo ativo de cada sessão.
Resumos dos prompts: auditar contratos, curar sem alterar origem, comparar
somente em treino/CV e reproduzir preservando modelos oficiais
(\texttt{USO\_DE\_IA.md}).

Decisões delegadas à IA receberam conferência humana posterior registrada
separadamente; o teste humano não foi gerado por IA.
Das 125 frases complementares atribuídas ao ChatGPT/GPT-6, 124 foram aproveitadas
e uma semelhante descartada. Outra candidata completou o treino com origem LLM
e revisão humana confirmada. Prompts recuperados estão em
\texttt{dados/prompts\_llm.md}; ferramenta e prompts de 158 frases permanecem
incompletos, sem reconstrução fictícia.

Jonathan usou Claude/claude.ai, Opus 5.5 conforme registro, inclusive na auditoria
que identificou a falha de reprodução; não utilizou Claude Code.
Felippe e Gabriel utilizaram inteligência artificial como apoio na elaboração do formulário, ajudando a reformular perguntas, melhorar a clareza das respostas e sugerir novas questões e ajustes para aprimorá-lo.
GPT-6 Luna/OpenRouter é o LLM do chatbot, distinto das ferramentas de desenvolvimento.

O Codex foi utilizado como ferramenta de apoio na organização de evidências, consulta à bibliografia, elaboração de tabelas e figuras e análise dos erros. Os prompts foram direcionados ao apoio na verificação das informações e à incorporação dos resultados da avaliação realizada, sem repeti-la. A avaliação foi feita manualmente pelo grupo.
Não houve chamadas ao OpenRouter ou acesso a credenciais nessas intervenções.

A IA acelerou organização e implementação, mas sugeriu guia desatualizado,
duplicações e confusão entre relato e norma. Foi conferida fontes,
confirmou a curadoria e preservou erros e exclusões. Os autores respondem por compreender e revisar cada resultado.
Análises assistidas não são apresentadas como independentes de IA.




\bibliographystyle{sbc}
\bibliography{referencias}
\end{document}
